\section{Moving insertions through a sewn surface}\label{subsec:motion}

The planting construction expresses the original tables on pairs of
spheres in each cell. We now combine these tables at a small number
of marks. Two outcomes of a Laurent expansion suffice. At an insertion
mark, a coefficient fuses the moving table with the table already
there. At a seam, a sufficiently negative coefficient transfers the
moving table to the neighboring cell and lowers the energy of the
seam. This decrease bounds both the number of crossings and the amount
of energy that any cell can receive.

In each cell, designate the marks of the original immobile stations
and one additional vacuum mark as \emph{collectors}. They stay in their
cell. Process the original movable stations in turn, skipping those
already absorbed. The moving table keeps its ordinary bit ports. If its
current legs on the sewn sphere $A$ and private sphere $B$ have
energies $m_A,m_B$, put
\[
       m=m_A+m_B/P,\qquad \mu=m+H,\qquad L=\lceil 2P\mu\rceil .
\]
Here $P$ is the exchange rate of Equation~\eqref{eq:plant-scales}.
The reserve $H$ stays with this move from its start until it ends, while
$m$ is updated after every fusion or crossing.

\subsection{Reconstructing an evaluation from collisions}

Fix the branch of the outer sum and, for the following identity, the
private selectors. The branch determines the energies, so the same
identity works for all allowed selector values. Hold all other positions
and insertion tables fixed. Each entry of the contraction is a rational
function of the mover's $B$-coordinate $z$. Its $A$-coordinate is $z^P$.
The sphere kernels have poles only at other marks on the current sphere
and are bounded at infinity. Near a mark with $B$-coordinate $z_0$, write
$s=z-z_0$. Since $P$ is a power of the characteristic, the corresponding
coordinate difference on $A$ is $s^P$. Hence an $A$-fusion of order $i$
and a $B$-fusion of order $v$ contribute to order $j=Pi+v$ of the
contraction. At an insertion of weighted energy $m_0$ the merged energy is
\begin{equation}
            m_{\rm new}=m_0+m+j/P.                     \label{eq:motion-fusion}
\end{equation}
Indeed each fusion adds its order to the two incident energies, by
Proposition~\ref{prop:fusion}. Nonnegativity of the two output energies
bounds $i$ and $v$ below. Thus the pole orders are bounded, and only
finitely many pairs $(i,v)$ contribute at each $j$.

To reconstruct the value at the current position, use the following
coefficients at the other marks of the cell.
\begin{itemize}
\item at every insertion, its entire principal part;
\item at the designated vacuum collector, also the orders $0,\ldots,gL$,
even after this collector has absorbed a table;
\item at each seam, only the orders $j\le -L$ within its pole bound.
\end{itemize}
A function on which all these coefficients vanish has poles only at the
at most $g$ seams, of order at most $L-1$ each. It has a zero of order
at least $gL+1$ at the designated collector and no pole at infinity.
The degree of the divisor of a nonzero rational function on the projective
line is zero, so this function vanishes identically. On the space with
the prescribed pole bounds, evaluation at the mover's position is
therefore a linear combination of the listed coefficients. This
combination acts entrywise, independently of the tables.

A coefficient at an insertion is a sum over the compatible pairs
$(i,v)$. For each pair, apply the two fusion maps and contract them
with the moving and stationary tables, including the sign of their
chosen order. A collision with a collector ends the move. At any
other insertion only $j<0$ is used, so the combined table continues
with at most the sum of the two original weighted energies. This step
can absorb an auxiliary insertion, an original station awaiting its
turn, or a private insertion left by an earlier crossing.

\subsection{Crossing a seam}

Suppose the selected coefficient has order $j\le-L$ at a seam of
energy $e$. Only sphere $A$ is sewn there. The insertion on $B$ at the
seam is the vacuum, so the $B$-expansion has Taylor orders $v\ge0$.
For each such $v$ with $i=(j-v)/P$ integral and in the fusion range,
the new $A$-energy at the seam satisfies
\begin{equation}
       e'=e+m_A+i\le e+m_A-2\mu\le e-\mu .             \label{eq:motion-drop}
\end{equation}
Terms with $e'<0$ vanish. With $i$ fixed, interpolate the $B$ Taylor
coefficient using evaluations at fresh regular positions. Each resulting
term has one private insertion of energy $m_B$ left behind, with
vacuum on $A$. Its index is carried to the original table by an ordinary
equality wire. We call this private insertion the \emph{residue} of the
crossing.

To identify the new moving table, fix its old $A$-state $x$ by
linearity. Let $a,b$ range over wedge basis states of energies $e,e'$,
respectively, and denote the two remaining sphere forms by
$C_{\rm near},C_{\rm far}$, suppressing their spectator arguments.
Lemma~\ref{lem:seam} gives the identity
\[
\begin{aligned}
 &\sum_b C_{\rm near}(b)\sum_a [I_{e'}(a,x)]_b\,C_{\rm far}(a)\\
 &\hspace{1em}=\alpha\sum_{t\le m_A}\sum_b C_{\rm near}(b)
                      \sum_a [I_e(b,S_t x)]_a\,C_{\rm far}(a).
\end{aligned}
\]
On the right, the inner sum is the Laurent coefficient of order
$e-e'-t$ for an insertion $S_t x$ approaching $b$ on the far-side
$A$-sphere. This order is integral, since $e'-e-m_A=i$ and $m_A-t$
are integers for nonzero terms. Interpolate the coefficient by far-side
regular evaluations, with vacuum in the corresponding $B$-position.
The sum over $b$ sews the two spheres at energy $e'$. The new moving
table, now on the far side with energies $(t,0)$, retains its ordinary
ports and the wire to the residue (Figure~\ref{fig:motion-crossing}).
All reordering signs depend only on the specified energies, and the
identity remains valid when $x$ is coupled to the private index or to a
spectator.

\begin{figure}[!htbp]
\centering
\begin{tikzpicture}[x=1cm,y=1cm, font=\small,
 cell/.style={draw=ink!50,fill=ink!3,rounded corners=2pt},
 tab/.style={draw=ink,fill=white,rounded corners=1pt,inner sep=3.2pt,
             font=\small,align=center}]
 \node[font=\small\bfseries] at (-3.62,1.45) {Near cell};
 \node[font=\small\bfseries] at (3.62,1.45) {Far cell};
 \draw[cell] (-5.36,-1.11) rectangle (-1.88,1.13);
 \draw[cell] ( 1.88,-1.11) rectangle ( 5.36,1.13);
 \node[align=center] at (-3.62,.55) {sewn sphere $A$};
 \node[align=center] at ( 3.62,.55) {sewn sphere $A$};
 \fill[sea] (-1.88,.55) circle (2pt);
 \fill[sea] ( 1.88,.55) circle (2pt);
 \draw[sea,line width=.9pt] (-1.88,.55) -- (1.88,.55);
 \node[fill=white,inner sep=2pt,text=sea] at (0,.55) {seam energy $e'$};
 \node[tab] (res) at (-3.62,-.40) {private residue\\ $(0,m_B)$};
 \node[tab] (mov) at ( 3.62,-.40) {moving table\\ $(t,0)$};
 \draw[dashed,rust,line width=.75pt] (res.east)--(mov.west);
 \node[fill=white,inner sep=2pt,text=rust] at (0,-.40) {$B$-state index};
\end{tikzpicture}
\caption{A moving insertion of energies $(m_A,m_B)$ crosses from the near
cell to the far cell. The solid wire identifies seam states of energy $e'\le e-\mu$. The dashed
wire carries the old private state to the new moving table, whose sewn
energy is $t\le m_A$.}
\label{fig:motion-crossing}
\end{figure}

Both interpolations can be made over the same finite field $K$ as the
original contraction. Multiplying by a common denominator embeds the
space of rational functions bounded at infinity with prescribed poles
into the polynomials of degree at most the total pole bound. Evaluation
at more than that many regular field points is injective, so these
evaluations span the dual and recover each required Laurent coefficient.
Exclude all occupied positions. The map $z\mapsto z^P$ permutes $K$, so
regular evaluations can use corresponding positions on both spheres.
The accounting below bounds the degrees by $O_D(T)$ and the number of
occupied positions per cell by a polynomial in $D$. One can first
obtain those bounds over the algebraic closure and use them to choose
$K$ before any interpolation. Fix an ordering of the field
elements and, for each allowed set of regular points in this order,
store the entries of one reconstruction matrix in the table dictionary. The local control need only
name the selected point and check its availability.

Allow at most $U=D^5$ crossings in a move. When this limit is reached,
its current mark becomes a collector. Nonterminal fusions between
successive crossings remove one stationary insertion each. A completed
move never releases a new mover, so this procedure terminates in every
summand.

\subsection{Accounting for motion}

We first bound the energy stored in the final insertions. An original
station of arity $t_i$ has initial weighted energy $O((1+t_i/T)^2H)$.
The auxiliaries contribute at most $DH$ per occupied cell. Also allocate
one reserve $H$ to each original movable station. The initial insertion
energy and reserves are $O(D^2H)$ per cell. Over all cells, including
all networks in the batch, the bound is
\[
 O\!\left(H\left(D\,\#\{\text{occupied cells}\}
                  +\#\{i\}+T^{-2}\sum_i t_i^2\right)\right)
       =O(WH/T).
\]

A continuing fusion cannot increase the sum of participating energies,
by Equation~\eqref{eq:motion-fusion}. Keep only the reserve of the
current move. In a crossing the energy $m$ pays for $t\le m_A$ on the
far side and $m_B/P$ at the near-side residue. At a final fusion the
collector may gain up to $m+gL/P\le(1+3g)\mu$, since $H\ge1$ and
$P\ge1$. Count energy at collectors with weight $1/(1+3g)$ and energy
elsewhere, including reserves, with weight one. The sum cannot
increase at any of the three operations or when a move expires. The
total insertion energy in every resulting summand is consequently
$O_g(WH/T)$.

There is also a bound in each cell. A crossing brings in at most
$t+H\le\mu$ of charge. By Equation~\eqref{eq:motion-drop} the seam loses
at least this much energy. A cell has at most $g$ seams, each initially
of energy $O_g(D^2H)$. Adding the initial charge within the cell bounds
its current insertion energy by $O_g(D^2H)$ at every step. Moreover,
each crossing decreases the seam energy by at least $H$, so throughout
the construction a cell has $O_g(D^2)$ crossing incidences. Only these
incidences can introduce new stationary residues. With $H=O_D(T)$ and
$P$ fixed, these energy and mark bounds give the asserted $O_D(T)$ total
pole bound. The number of expired moves over all cells is at most the
total number of crossings divided by $U$. Including the initial immobiles
and designated collectors, the total number of collectors is $O(W/(DT))$.

It remains to route the ordinary wires created in the expansion. By
Equation~\eqref{eq:plant-widths}, the private state left at a crossing
uses $O_g(DT/\sqrt F+1)$ bits. One endpoint follows the mover. The
other stays at the residue until another insertion absorbs it, then
travels with that insertion. A residue that can still be absorbed
carries only the endpoint created at its own crossing. Once a moving
endpoint finishes a move, its collector stays fixed. A move visits at
most $U+1$ cells counting repetitions, collecting at most $O_g(D^2)$
residue endpoints per visit and creating at most $U$ of its own.
Hence it carries $O_g((U+1)D^2)$ endpoints along any one hop. Using
the bound on incident crossings and including terminations gives the
looser bound $O_g((U+1)D^6)$ local wire segments per cell. Successive
motions within a single cell require only one local connection. With
$F$ from Section~\ref{subsec:planting}, these segments have total width
$O(D^{-2\ell-1}T)+O_D(1)$ there.

The private-selector wires start at marks that read them and at the
stations checking the selectors and evaluating the sphere forms. Their
endpoints follow the same paths. Each cell starts with polynomially
many such ports in $D$ and the number of chunks, so the path and
crossing bounds give a polynomial bound in the same quantities on
incident selector segments and copies. This count is independent of
$b$ and of the selector rate. After fixing the number of chunks, choose
the rate small enough that the total selector width, including equality
copies, is at most $D^{-2\ell}T$ per cell. The next section converts the
remaining energy levels into table arities.
